\section{Conclusão}

\paragraph{}
As atividades permitiram compreender na prática como ataques de engenharia
social podem ser utilizados para induzir usuários a fornecer informações
sensíveis ou executar ações que comprometam a segurança da organização.

\paragraph{}
A utilização das ferramentas SET, Gophish e BeEF demonstrou diferentes
abordagens para execução de campanhas de phishing e coleta de informações em
ambientes controlados, evidenciando a importância do fator humano na segurança
da informação.

\paragraph{}
Os resultados reforçam a necessidade de programas contínuos de conscientização,
autenticação multifator, monitoramento de acessos e políticas de segurança bem
definidas para reduzir a eficácia desse tipo de ataque.
